\section{Attainment, classical inputs, and realization theory}
\label{sec:v33-comparison}

\subsection{Clustered spectra and attained prediction laws}
\label{sec:v35-spectrum}
The full spectrum, rather than only its least singular value, is a
natural comparison object. Write
\[
 F_L(x)=(e^{\mathrm i kx_j})_{0\le k\le L,\,1\le j\le q}.
\]
Batenkov, Diederichs, Goldman and Yomdin
\cite[Theorems 2.2--2.3]{BDGGY2021} study this Fourier--Vandermonde
matrix on the circle. For an $s$-point cluster with all pairwise
circular distances between $\tau h$ and $h$, their single-cluster
theorem gives
\[
 \sigma_j(F_L)\asymp L^{1/2}(Lh)^{j-1},\qquad 1\le j\le s,
\]
when $L\ge s$ and $Lh$ is sufficiently small, with constants depending
on $s,\tau$. Their multi-cluster theorem compares the complete spectrum
with the ordered union of the cluster spectra, with multiplicative
errors controlled by $(L\vartheta)^{-1}+Lh$, where $\vartheta$ is
inter-cluster separation and $h$ the largest cluster diameter, in the
stated range $C/\vartheta\le L\le c/h$. Thus the spectral comparison
is not confined to the worst-conditioned direction.

For our fixed horizon, the connection to those matrix objects admits
a short direct proof. It includes exact repeated labels and requires
no choice of a collision path.

\begin{proposition}[Exterior spectral products and determinant volumes]
\label{prop:v35-exterior-spectrum}
Fix integers $q\ge1$ and $L\ge q-1$. For $x\in[0,1]^q$, with
repeated entries permitted, let $V_\ell(x)$ be the largest absolute
Vandermonde product on $\ell$ entries, with $V_1=1$. If
$\sigma_1\ge\cdots\ge\sigma_q\ge0$ are the singular values of
$F_L(x)$, there are constants $c_q,C_{q,L}>0$ such that
\begin{equation}\label{eq:v35-spectral-volume}
 c_q V_\ell(x)\le\prod_{j=1}^{\ell}\sigma_j
          \le C_{q,L}V_\ell(x),\qquad 1\le\ell\le q.
\end{equation}
In particular both sides vanish exactly when fewer than $\ell$
distinct node values remain. The constants are independent of all
node gaps; $q$ and $L$ are fixed.
\end{proposition}
\begin{proof}
Put $z_j=e^{\mathrm i x_j}$. For $|x_j-x_i|\le1$,
\[
 2\sin(1/2)|x_j-x_i|\le |z_j-z_i|\le |x_j-x_i|.
\]
A minor with columns $j_1,\ldots,j_\ell$ and rows
$0\le k_1<\cdots<k_\ell\le L$ is an alternating polynomial
$\det(z_{j_b}^{k_a})_{a,b=1}^{\ell}$. It is divisible by every
$z_{j_b}-z_{j_a}$, and hence by their product: these distinct linear
factors are relatively prime in the polynomial ring. The quotient
is a polynomial. There are finitely many row choices, so their
quotients have a common bound on the unit torus, depending only on
$\ell,L$. Every $\ell$-minor is therefore bounded by a constant
times $V_\ell(x)$. This polynomial factorization remains valid
at repeated nodes without division by a numerical gap.

The matrix of $\bigwedge^\ell F_L$ in the standard exterior bases
consists of these minors. Its operator norm is bounded by its
Frobenius norm, hence by
$\binom{L+1}{\ell}^{1/2}\binom q\ell^{1/2}$ times their largest
absolute value. Conversely, take columns attaining $V_\ell$ and
the first $\ell$ rows. Their minor is precisely the Vandermonde
product in the corresponding $z$'s, so its absolute value is at
least $(2\sin(1/2))^{\ell(\ell-1)/2}V_\ell(x)$.
The operator norm dominates any matrix entry. Finally, the singular
value decomposition gives
$\|\bigwedge^\ell F_L\|=\prod_{j\le\ell}\sigma_j$.
Taking the minimum and maximum of the finitely many constants proves
\eqref{eq:v35-spectral-volume}, including all zero-rank cases.
\end{proof}

This proposition uses elementary exterior algebra and polynomial
factorization; it is a comparison identity, not a new clustered-matrix
estimate with uniform growing-bandwidth constants. In particular it
does not replace the dependence on $L$ established in the cited work.

\begin{corollary}[Spectral expression of the acquired envelope]
\label{cor:v35-spectral-envelope}
For each $1\le m<N$, fix an integer $L_m\ge q_m-1$ and apply
$F_{L_m}$ to the formal positive nodes $x_\alpha(a)$ of
Section~\ref{sec:collision-geometry}. Write its singular values as
$\sigma_{m,1}(a)\ge\cdots\ge\sigma_{m,q_m}(a)$. Then
\begin{equation}\label{eq:v35-spectral-envelope}
 \Psi_{n,m}(M,a)\asymp
 \max_{1\le\ell\le p_{n,m}}
    \left(\frac{\prod_{j=1}^{\ell}\sigma_{m,j}(a)}{M}\right)^{2/\ell}.
\end{equation}
The comparison is uniform in every integer $M\ge1$ and $a\in K$,
including exact collisions. The same spectral envelope, maximized
over $n+m=N$, is comparable to either common-controller risk in
Theorem~\ref{thm:resolution-main}.
\end{corollary}
\begin{proof}
Use Proposition~\ref{prop:v35-exterior-spectrum} with
$q=q_m$ and $V_\ell=\mathcal V_{m,\ell}$. Raise its inequalities
to $2/\ell$ and maximize over the finite set of allowed indices.
The last assertion is Theorem~\ref{thm:resolution-main}.
\end{proof}

The matrix in this corollary is an auxiliary function of the known
exponent labels, not an extra observation supplied to the predictor.
Its full spectrum describes the available collision scales. The
upper endpoint $p_{n,m}=\min\{n(r-1),q_m\}$ and the interpretation
of those scales as prediction regret require the acquired tangent,
probability minorization, whole-image cover and causal update proved
in the preceding sections. The following example separates these
roles while holding the future spectrum fixed.

\begin{corollary}[Dependence on acquisition at a fixed future test space]
\label{cor:v35-fixed-future}
Use the detector, prior and command cube of
Corollary~\ref{cor:two-parameter}, and keep the same two-trial query
menu. For any fixed $L\ge8$, the nine singular values of the
auxiliary matrix on its positive future nodes satisfy, componentwise,
\begin{equation}\label{eq:v35-two-parameter-spectrum}
 (\sigma_{2,1},\ldots,\sigma_{2,9})
       \asymp(1,1,1,1,1,1,\rho,\rho,\tau)
\end{equation}
uniformly on the calibration square. The checkpoint problems
$(n,2)$, $n=1,2,3$, with their respective total horizons $n+2$,
have unconditional optimal regrets
\begin{equation}\label{eq:v35-fixed-future}
 \begin{split}
 \overline R_{M;1,2}^{u,v}&\asymp M^{-2/3},\\
 \overline R_{M;2,2}^{u,v}&\asymp M^{-1/3},\\
 \overline R_{M;3,2}^{u,v}&\asymp
  \max\{M^{-1/3},\rho^{1/2}M^{-1/4},
                       (\rho^2\tau)^{2/9}M^{-2/9}\}.
 \end{split}
\end{equation}
The same comparisons hold for their worst-history checkpoint optima.
\end{corollary}
\begin{proof}
The six separated groups in the proof of
Corollary~\ref{cor:two-parameter} give
$\mathcal V_{2,\ell}\asymp1$ for $\ell\le6$ and successive
volumes of orders $\rho,\rho^2,\rho^2\tau$ for $\ell=7,8,9$.
Proposition~\ref{prop:v35-exterior-spectrum} gives the corresponding
partial spectral products. Dividing successive positive products
gives \eqref{eq:v35-two-parameter-spectrum}. At $\tau=0$ with
$\rho>0$ the last singular value is zero, and at $\rho=0$ the last
three vanish, by the exact-rank clause of the proposition.
The truncations are $p_{1,2}=3$, $p_{2,2}=6$, $p_{3,2}=9$.
Substituting them in Theorem~\ref{thm:intrinsic-checkpoint} gives
the first two laws; the third is the checkpoint computation in
Corollary~\ref{cor:two-parameter}.
\end{proof}

Thus one and the same future spectrum is compatible with different
acquired dimensions and resolution envelopes. This is a consequence
of the classification, not a claim that singular values contain no
useful collision information. Nor do the three checkpoint problems
assert a new sharp price for a multi-checkpoint controller; their
horizons are stated separately.

\subsection{Superresolution minimax loss and retained-label loss}
Batenkov, Demanet, Goldman and Yomdin
\cite[Definition 3.13 and Theorem 3.14]{BDGY2020} obtain statistical
minimax consequences, not only conditioning estimates. They reconstruct
complex coefficients of $s$ point sources on a grid of spacing
$\Delta$ from the entire noisy Fourier function on $[-\Omega,\Omega]$.
The support is also unknown. For normalized $L^2$ noise at most
$\delta$, the coefficient-$\ell^2$
minimax error has upper and lower estimates on the scale
$\delta\,\mathrm{SRF}^{2b-1}$, with
$\mathrm{SRF}=\pi/(\Delta\Omega)$ and clustering parameter $2\le b\le s$.
The upper statement selects arbitrarily large fixed values of
$\mathrm{SRF}$ and then sufficiently small $\Delta$; the lower
statement uses specified cluster parameters and all sufficiently
large $\mathrm{SRF}$ in its stated grid range. These are the separate
quantifiers of that theorem, not a bound asserted for every bandwidth
and configuration.

Here acquisition is a fixed stochastic experiment, and the latent
parameter is not an unknown coefficient vector on a reconstruction
grid. The loss compares a retained prediction with the full-history
prediction in the same experiment. The integer $M$ constrains the
number of labels between reports; it is neither a Fourier sample
count nor the reciprocal of external observation noise. Even the
worst-history criterion concerns prediction after compression, not
recovery of a sparse signal from perturbed spectral data. Consequently
one cannot replace $\Delta\Omega$ or $\delta$ by $M^{-1}$ in the
superresolution formula to obtain our theorem. The relevant common
mathematics is degeneration of exponential test systems; the conversion
to the particular loss and the attained family is different.

\begin{remark}[Determinant laws beyond finite-power paths]
\label{rem:v35-flat-path}
Theorem~\ref{thm:collision-tree} has an additional hypothesis comparing
each nonidentical pair gap to a finite power of its path parameter.
Neither Theorem~\ref{thm:resolution-main} nor
Corollary~\ref{cor:v35-spectral-envelope} uses that hypothesis.
For example set, for sufficiently small $\theta>0$,
\[
 u=e^{-1/\theta^2},\qquad
 v=e^{-1/\theta^2}+e^{-1/\theta^4},
\]
and extend by $u=v=0$ at zero. Then
$\rho\asymp e^{-1/\theta^2}$ and
$\tau=e^{-1/\theta^4}$. Direct substitution in
\eqref{eq:two-parameter-law} gives the three resolution regimes;
the two crossover budgets have orders
\[
 \rho^{-6}\asymp e^{6/\theta^2},\qquad
 \rho^2\tau^{-8}\asymp e^{8/\theta^4-2/\theta^2}.
\]
They follow by balancing adjacent terms. The relevant gaps are not
comparable to any finite powers of $\theta$, so this example uses
the determinant theorem directly rather than its finite-power tree
specialization.
\end{remark}


\subsection{What the geometric classification requires}
The distinction between a test space and its attained prediction image
is essential to the contribution of Theorem~\ref{thm:resolution-main}.
Generalized Vandermonde determinants, confluent divided differences and
Leja ordering are classical \cite{Pinkus,deBoor,Leja,BosEtAl}. For a
supplied exponent configuration they describe a system of coordinates
and its degenerating scales. They do not identify the tangent of the
likelihood products that the command cube can produce, or show that
the induced posterior law has mass in each relevant direction.
Lemma~\ref{lem:binomial-tangent} supplies that attainable tangent.
Its pairing with complete confluent flags, followed by normalization,
is the experiment-specific step in Lemma~\ref{lem:newton-attainment}.

Likewise, a volume argument for quantizing a supplied rectangle gives
a lower bound only after the rectangle has been obtained under the
law used in the risk. Here the minorization is for the actual joint
command-and-report law, with the all-failure evidence retained.
The full-support prior may be singular: the local absolutely continuous
variables are the acquisition commands. The dimension truncation is
therefore supported by an acquired differential and a positive-mass
chart, not by an ambient rank count. The normalized product loses
exactly one rank because the unnormalized product belongs to its tangent.

For the converse direction, the real entropy inequality recalled in
\cite{YomdinComte,ComteHalupczok} and the bounded-format section bound
in \cite[Lemma 2.18]{ZhangKileel} are classical inputs. The application
here first places the entire attained image, over every report word,
in a bounded-format thin rectangle. Its format bound is independent
of the real moment and exponent coefficients. The resulting covering
products stop at the dimension of the image. A local positive-mass
chart, by itself, would not justify this global upper bound.

Finally, quantization followed by a stable filtering error recurrence
is not a new general principle; the nonlinear-filter comparisons are
given in the companion, Section~\ref{sec:filter-quantization-comparison}.
A signal grid with finitely many sites may still retain a continuous
vector of posterior weights. The resource here is instead a finite
number of possible retained posterior representatives. The raw update
in \eqref{eq:formal-multiindex-update} maps reachable representatives
to reachable representatives and has a gap-independent Lipschitz bound.
It makes the global covers compatible with one causal label process.
The same acquired law supplies matching lower bounds at all checkpoints.
Together these steps identify the full all-budget collision profile;
none follows from an interpolation estimate or a filtering upper bound
alone. These comparisons identify the role of the cited inputs, not
an exhaustive priority theorem.

\subsection{The Blackwell region and the finite-cell algebra}
Comparison of experiments is the immediate framework for the one-step
coefficient set \cite{Blackwell1953}. Proposition~\ref{prop:v34-blackwell}
gives the exact identity $\mathscr A_M(\nu)=\{\mathsf E U\}$, not
merely an analogy. The support formula is its finite-action Bayes reward.
The row-overlap inequality expresses preservation of a common
submeasure under postprocessing. Neither fact is presented as a new
general comparison-of-experiments theorem. Their acceptance/failure
form makes explicit which coefficients a physical transition can realize.
In particular, the auxiliary raw index is not revealed after failure.

Iterating those coefficients gives the conditional occupancy polynomials.
Theorem~\ref{thm:v33-moment-program} states the resulting exact
mean-risk formula on a general standard Borel latent space with finite
nonnegative cells, including zero evidence. The compact coefficient
representation removes history-indexed optimization variables and
characterizes one common controller. It remains a nonconvex optimization;
it does not by itself identify a sharp causal price, an optimal label
allocation between checkpoints, or an optimizer phase diagram.
The evaluation identity in Proposition~\ref{prop:v34-evaluation}
connects this general algebra to $W_{mA}$ but does not provide the
monomial tangent or its attained collision scales. Those are separate
conclusions of the first part of the article.

The finite-action purification mechanism is classical
\cite{DWW1951}; its proof is included to specify precisely which
integrals are preserved. Here they are all the implementability
coefficients, including the common failure coefficients. The result
preserves conditional state laws and the entire vector of mean regrets
under atomless acquisition. It neither purifies individual-history
risk nor removes a persistent public seed. Polyhedral realization
concerns scalarized mean risks, with atomic ties retained, and entails
no minimax exchange. Compact collision transfer is a fixed-$M$
continuity statement. The all-budget scale of that same exact formula
comes from Theorem~\ref{thm:resolution-main}, not from compactness.

\subsection{Finite approximation and information resources}
The finite-approximation principles used here have several distinct
antecedents. We specify their relation to the retained resource, rather
than identify finiteness of a model with finiteness of its causal memory.

Saldi, Y\"uksel and Linder \cite[Assumption 3.1 and Theorem 3.2]{SYLWeak2016}
prove compact-uniform convergence of discounted values when an action
space is replaced by increasingly fine finite nets. Their assumptions
include continuity of the cost, weak continuity of the controlled kernel,
compactness of the action space, and a weight condition. Their feedback
policies use the model state; the partially observed application uses
a belief state. Our resource is instead the number of labels that may
persist between reports. Neither the full history nor its belief state
is an available feedback argument. Action quantization and persistent
state compression are therefore different approximations. The
fixed-horizon bound in Theorem~\ref{thm:v32-certificates} also specifies
a finite global enclosure and its error terms; it is not an application
of a compact-uniform value limit with an unspecified rate.

For decentralized control, Saldi, Y\"uksel and Linder
\cite[Theorems 4, 5 and 11]{SYLTeams2015} give finite observation and
action approximations and asymptotically optimal lifted policies under
the stated continuity, compactness or integrability hypotheses. The
information available to each decision maker is part of their team
model. The distinction here is not an absence of information constraints
in team theory. It is the particular constraint that later decisions
must factor through the same retained label. A representation giving a
new decision maker the entire acquired prefix would describe a larger
class unless that factorization were imposed. The compatibility
identities in Section~\ref{sec:v32-compatibility} impose it explicitly;
Theorem~\ref{thm:v33-moment-program} eliminates the prefix variables
for prescribed independent acquisition with a finite set of cells.

Y\"uksel and Linder \cite[Theorems 3.4 and 6.2]{YLChannels2012}
establish continuity of optimal costs under total-variation convergence
of observation channels, including a finite-horizon result for uniform
convergence in the state under bounded measurable costs. Their policies
may use their observation and action histories. The lifted-command
argument here similarly compares laws on a common observation space,
so that no regularity of an encoder's decision regions is needed. It
first compares the true experiment with a cellwise-constant likelihood
on that same space. Merely replacing a continuous command law by point
masses would not establish total-variation convergence. After this
comparison, compatibility restricts both problems to the same label
resource. Calibration uniformity comes from the positive likelihood
margin; the collision-dependent risk comparison additionally requires
the acquired geometry proved earlier.

Finite-controller optimization and finite-action purification also
precede the present work; see \cite{V32Amato,DWW1951}. The matrix body
has the explicit support formula \eqref{eq:v33-support} and compulsory
row overlap in Corollary~\ref{cor:v33-overlap}; replacing it by arbitrary
stochastic rows changes the model. Its polynomial iteration, deterministic
and scalarized polyhedral realizations, and compact collision limits
are a realization theory for the same per-report resource. None requires
regeneration or free persistent side information. The graph and
regenerative developments in the companion retain their separate
hypotheses and complete proofs.
